% ==========================================
% LaTeX-Vorlage: Hausarbeit
% Thema: Bestellanforderungsprozess mit SAP BTP
% ==========================================

\input{settings}
\makeglossaries

\begin{document}

% Titelseite
\newgeometry{left=2cm,right=2cm,top=2cm,bottom=2cm}
\begin{titlepage}
    \centering
    \Large

    {\LARGE \textbf{Konzeption und Implementierung eines Bestellanforderungsprozesses mit Genehmigungsworkflow auf der SAP Business Technology Platform}}

    \vfill
    Hausarbeit \\
    vorgelegt von \\
    Dominik Sitny \\

    \vspace*{2cm}
    Angefertigt im praxisintegrierten Bachelorstudiengang \\
    Wirtschaftsinformatik -- Bachelor of Science (B.Sc.)

    \vspace*{1cm}
    am Fachbereich Wirtschaft \\
    der Hochschule Bielefeld \\
    Wintersemester 2025/26

    \vspace*{5cm}
    {\normalsize
    Erstprüfer/in: \hspace{0.5cm} Prof.\ Dr.\ N.\,N. \\
    Zweitprüfer/in: \hspace{0.5cm} Prof.\ Dr.\ N.\,N. \\
    }

\end{titlepage}

\restoregeometry

% Abstract
\cleardoublepage
\thispagestyle{empty}
\section*{Abstract}

Die vorliegende Hausarbeit untersucht die Konzeption und prototypische Implementierung eines Bestellanforderungsprozesses mit integriertem Genehmigungsworkflow auf der SAP Business Technology Platform (BTP). Im Mittelpunkt steht die Frage, wie sich durch die Kombination des SAP Cloud Application Programming Model (CAP), SAP Fiori Elements und SAP Build Process Automation ein durchgängiger, cloudnativer Geschäftsprozess realisieren lässt. Auf Basis einer Anforderungsanalyse wird zunächst ein Architekturkonzept entworfen, das eine klare Schichtentrennung zwischen Datenmodell, Geschäftslogik und Benutzeroberfläche vorsieht. Anschließend wird der Prototyp implementiert, der es Fachanwendern ermöglicht, Bestellanforderungen mit Dokumentenanhängen in einer Fiori-Elements-Oberfläche zu erstellen und zur Genehmigung einzureichen. Der Genehmigungsworkflow wird in SAP Build Process Automation modelliert und über eine REST-Schnittstelle bidirektional mit dem CAP-Service gekoppelt. Der Workflow-Kontext umfasst Positionslisten, Anhänge-Metadaten mit direkten Download-URLs sowie Deep Links zur Fiori-Anwendung. Eine Statushistorie protokolliert den vollständigen Lebenszyklus jeder Bestellanforderung. Die Ergebnisse zeigen, dass die gewählte Architektur eine effiziente Entwicklung ermöglicht und die Plattformdienste der SAP BTP eine nahtlose Integration von Anwendungslogik und Prozessautomatisierung erlauben.

\bigskip
\large{\textbf{Abstract (English)}}

\normalsize
This paper examines the design and prototypical implementation of a purchase requisition process with an integrated approval workflow on the SAP Business Technology Platform (BTP). The central question is how a cloud-native end-to-end business process can be realized by combining the SAP Cloud Application Programming Model (CAP), SAP Fiori Elements, and SAP Build Process Automation. Based on a requirements analysis, an architecture concept is developed that provides a clear separation of concerns between data model, business logic, and user interface. Subsequently, a prototype is implemented that enables business users to create purchase requisitions with document attachments in a Fiori Elements interface and submit them for approval. The approval workflow is modeled in SAP Build Process Automation and bidirectionally coupled with the CAP service via a REST interface. The workflow context includes item lists, attachment metadata with direct download URLs, and deep links to the Fiori application. A status history records the complete lifecycle of each purchase requisition. The results demonstrate that the chosen architecture enables efficient development and that the platform services of SAP BTP allow seamless integration of application logic and process automation.


\pagenumbering{Roman}

% Inhaltsverzeichnis
\cleardoublepage
\phantomsection
\tableofcontents
\cleardoublepage

% Tabellenverzeichnis
\phantomsection
\addcontentsline{toc}{section}{Tabellenverzeichnis}
\listoftables
\cleardoublepage

% Abkürzungsverzeichnis
\phantomsection
\addcontentsline{toc}{section}{Abkürzungsverzeichnis}
\printglossary[type=\acronymtype,title={Abkürzungsverzeichnis},nonumberlist]
\newacronym{ki}{KI}{Künstliche Intelligenz}
\newacronym{btp}{BTP}{Business Technology Platform}
\newacronym{cap}{CAP}{Cloud Application Programming Model}
\newacronym{cds}{CDS}{Core Data Services}
\newacronym{spa}{SPA}{SAP Build Process Automation}
\newacronym{mta}{MTA}{Multi-Target Application}
\newacronym{hdi}{HDI}{HANA Deployment Infrastructure}
\newacronym{odata}{OData}{Open Data Protocol}
\newacronym{cf}{CF}{Cloud Foundry}
\newacronym{xsuaa}{XSUAA}{Extended Services for User Account and Authentication}
\newacronym{ui5}{UI5}{SAPUI5 -- SAP User Interface Development Toolkit for HTML5}
\newacronym{crud}{CRUD}{Create, Read, Update, Delete}
\newacronym{rest}{REST}{Representational State Transfer}
\newacronym{api}{API}{Application Programming Interface}
\newacronym{bpm}{BPM}{Business Process Management}
\newacronym{erp}{ERP}{Enterprise Resource Planning}
\newacronym{saas}{SaaS}{Software as a Service}
\newacronym{paas}{PaaS}{Platform as a Service}
\newacronym{jwt}{JWT}{JSON Web Token}


\newcounter{savepage}
\setcounter{savepage}{\value{page}}

\pagenumbering{arabic}

\section{Einleitung}

Unternehmen stehen zunehmend vor der Herausforderung, ihre internen Geschäftsprozesse zu digitalisieren und zu automatisieren, um wettbewerbsfähig zu bleiben. Ein zentraler Bestandteil betrieblicher Abläufe ist der Beschaffungsprozess, insbesondere die Bestellanforderung und deren Genehmigung. Die vorliegende Hausarbeit befasst sich mit der Konzeption und prototypischen Implementierung eines solchen Bestellanforderungsprozesses mit integriertem Genehmigungsworkflow auf der \gls{btp}.

\subsection{Motivation und Problemstellung}

In vielen Unternehmen stellt die Beschaffung einen geschäftskritischen Prozess dar, der unmittelbar die Kostenkontrolle, die Einhaltung von Budgetvorgaben und die Lieferfähigkeit beeinflusst \citep{weske2019bpm, dumas2018bpm}. Bestellanforderungen müssen in der Regel einen mehrstufigen Genehmigungsprozess durchlaufen, bevor sie in eine verbindliche Bestellung überführt werden. Dieser Genehmigungsprozess dient der Sicherstellung, dass Ausgaben im Rahmen der Unternehmensrichtlinien liegen und von den zuständigen Verantwortlichen freigegeben werden.

In der Praxis sind derartige Genehmigungsprozesse jedoch häufig noch durch manuelle, papierbasierte oder E-Mail-gestützte Abläufe gekennzeichnet. Dies führt zu einer Reihe von Problemen: Genehmigungsanfragen gehen verloren oder werden verzögert bearbeitet, die Nachvollziehbarkeit einzelner Entscheidungen ist eingeschränkt, und die Einhaltung von Compliance-Vorgaben lässt sich nur schwer sicherstellen. Darüber hinaus fehlt es an einer zentralen Übersicht über den aktuellen Status offener Bestellanforderungen, was die Transparenz im gesamten Beschaffungsprozess erheblich beeinträchtigt.

Vor dem Hintergrund der digitalen Transformation gewinnen integrierte Plattformlösungen an Bedeutung, die es ermöglichen, Geschäftsprozesse durchgängig zu modellieren, zu automatisieren und zu überwachen. Im SAP-Ökosystem bietet die \gls{btp} eine umfassende Plattform für die Entwicklung cloudbasierter Geschäftsanwendungen \citep{sap2024btp}. In Kombination mit dem \gls{cap} für die Anwendungsentwicklung und \gls{spa} für die Workflow-Automatisierung ergibt sich die Möglichkeit, Beschaffungsprozesse effizient und standardisiert abzubilden. Durch den Einsatz moderner \gls{bpm}-Konzepte können Unternehmen ihre Prozesse nicht nur digitalisieren, sondern auch kontinuierlich optimieren \citep{weske2019bpm}.

\subsection{Zielsetzung}

Ziel dieser Arbeit ist es, einen Prototyp für einen Bestellanforderungsprozess mit integriertem Genehmigungsworkflow auf der \gls{btp} zu konzipieren und zu implementieren. Dabei steht die folgende Forschungsfrage im Mittelpunkt: \textit{Wie lässt sich durch die Kombination von \gls{cap}, SAP Fiori Elements und \gls{spa} ein durchgängiger, cloudnativer Bestellanforderungsprozess mit Genehmigungsworkflow auf der \gls{btp} realisieren?} Die Anwendung soll es Mitarbeitenden ermöglichen, Bestellanforderungen über eine webbasierte Benutzeroberfläche zu erstellen, die anschließend automatisiert einem definierten Genehmigungsprozess zugeführt werden.

Für die technische Umsetzung wird das \gls{cap} mit Node.js als serverseitiges Programmiermodell verwendet \citep{sap2024cap}. Die Benutzeroberfläche wird auf Basis von SAP Fiori Elements generiert, um eine konsistente und standardkonforme Benutzererfahrung zu gewährleisten. Die Genehmigungslogik wird mithilfe von \gls{spa} als Workflow realisiert, sodass Genehmigungsschritte automatisiert ausgelöst und nachverfolgt werden können. Als Datenbank kommt SAP HANA Cloud zum Einsatz, und die Bereitstellung der Gesamtanwendung erfolgt über Cloud Foundry auf der \gls{btp}.

Im Einzelnen verfolgt die Arbeit die folgenden Teilziele:

\begin{itemize}
    \item Analyse der fachlichen Anforderungen an einen Bestellanforderungsprozess mit Genehmigungsworkflow,
    \item Konzeption einer geeigneten Anwendungsarchitektur unter Verwendung der genannten SAP-Technologien,
    \item Prototypische Implementierung der Anwendung einschließlich Datenmodell, Geschäftslogik, Benutzeroberfläche und Workflow,
    \item Bereitstellung der Anwendung auf der \gls{btp} und Evaluation des Ergebnisses.
\end{itemize}

\subsection{Aufbau der Arbeit}

Die vorliegende Arbeit gliedert sich in sechs Kapitel. Nach dieser Einleitung werden in \textbf{Kapitel~2} die theoretischen und technologischen Grundlagen erläutert. Dies umfasst eine Einführung in das \gls{bpm}, die \gls{btp} und deren relevante Dienste sowie das \gls{cap} und \gls{spa}. In \textbf{Kapitel~3} erfolgt die Konzeption der Anwendung. Hier werden die fachlichen Anforderungen analysiert, das Datenmodell entworfen und die Systemarchitektur beschrieben. \textbf{Kapitel~4} widmet sich der technischen Implementierung des Prototyps. Dabei wird auf die Umsetzung des Datenmodells, der Geschäftslogik, der Benutzeroberfläche mit SAP Fiori Elements sowie des Genehmigungsworkflows mit \gls{spa} eingegangen. In \textbf{Kapitel~5} werden die erzielten Ergebnisse vorgestellt und kritisch bewertet. Abschließend fasst \textbf{Kapitel~6} die wesentlichen Erkenntnisse der Arbeit zusammen und gibt einen Ausblick auf mögliche Erweiterungen.

\section{Grundlagen}
\label{sec:grundlagen}
Das folgende Kapitel legt die theoretischen und technischen Grundlagen dar, die für das Verständnis der in dieser Arbeit vorgestellten Lösung erforderlich sind. Zunächst wird die übergeordnete Plattform \gls{btp} eingeführt, auf der sämtliche Komponenten betrieben werden. Anschließend werden das \gls{cap}, die SAP-Fiori-Designrichtlinien, \gls{spa}, das \gls{odata}-Protokoll sowie das Konzept der Multi-Target Applications erläutert.

\subsection{SAP Business Technology Platform}
\label{subsec:btp}
Die SAP Business Technology Platform (\gls{btp}) ist die zentrale Cloud-Plattform der SAP SE, die Unternehmen eine integrierte Umgebung für die Entwicklung, Integration und Erweiterung betriebswirtschaftlicher Anwendungen bereitstellt \citep{metzger2021sap}. Sie vereint mehrere strategische Leistungsbereiche unter einem gemeinsamen Dach: Datenbank- und Datenmanagement, Analytik, Anwendungsentwicklung, Integration sowie Künstliche Intelligenz \citep{sap2024btp}.

Ein wesentliches Merkmal der \gls{btp} ist ihr Multi-Cloud-Ansatz. Die Plattform kann auf den Infrastrukturen der Hyperscaler Amazon Web Services, Microsoft Azure und Google Cloud Platform betrieben werden. Dadurch erhalten Unternehmen die Flexibilität, ihre Anwendungen in der jeweils bevorzugten Cloud-Umgebung bereitzustellen, ohne an einen einzelnen Infrastrukturanbieter gebunden zu sein.

Aus Sicht der Anwendungsentwicklung stellt die \gls{btp} verschiedene Laufzeitumgebungen zur Verfügung \citep{leymann2011cloud}. Die in dieser Arbeit verwendete Umgebung basiert auf \gls{cf}, einer quelloffenen \gls{paas}-Technologie. \gls{cf} ermöglicht es Entwicklern, Anwendungen in verschiedenen Programmiersprachen bereitzustellen, wobei die Plattform die zugrunde liegende Infrastruktur -- einschließlich Skalierung, Netzwerk und Betriebssystemverwaltung -- abstrahiert. Neben der Laufzeitumgebung bietet die \gls{btp} eine Vielzahl von verwalteten Diensten an, darunter Datenbanken wie SAP HANA Cloud, Authentifizierungsdienste (\gls{xsuaa}), Destination Services für die Anbindung externer Systeme sowie Dienste für die Prozessautomatisierung.

Die Administration der \gls{btp} erfolgt über das SAP BTP Cockpit, eine webbasierte Verwaltungsoberfläche. Hier werden Subaccounts, Umgebungen, Service-Instanzen und Berechtigungen konfiguriert. Subaccounts bilden dabei die organisatorische Einheit, in der Ressourcen und Berechtigungen gebündelt werden.

\subsection{SAP Cloud Application Programming Model}
\label{subsec:cap}
Das SAP Cloud Application Programming Model (\gls{cap}) ist ein quelloffenes Framework der SAP, das speziell für die Entwicklung von Unternehmensanwendungen auf der \gls{btp} konzipiert wurde. Es verfolgt einen deklarativen Ansatz, bei dem fachliche Domänenmodelle und Servicedefinitionen im Mittelpunkt stehen, während technische Aspekte wie Datenbankanbindung, Protokollhandhabung und Authentifizierung durch das Framework bereitgestellt werden \citep{sap2024cap}.

Das zentrale Beschreibungsmittel von \gls{cap} sind die Core Data Services (\gls{cds}) \citep{herger2020cap}. Mit \gls{cds} werden sowohl die Datenmodelle als auch die Serviceschnittstellen in einer kompakten, domänenspezifischen Sprache definiert. Aus einer \gls{cds}-Modellierung erzeugt das Framework automatisch die zugehörigen Datenbankartefakte, \gls{odata}-Serviceschnittstellen sowie die benötigte Laufzeitinfrastruktur. Dieser generative Ansatz reduziert den manuellen Implementierungsaufwand erheblich und stellt gleichzeitig die Einhaltung von Standards sicher.

\gls{cap} bietet darüber hinaus eine Reihe eingebauter Funktionalitäten, die für betriebswirtschaftliche Anwendungen typischerweise benötigt werden. Dazu zählen die automatische Bereitstellung von \gls{odata}-V4-Schnittstellen, die Integration mit dem \gls{xsuaa}-Dienst zur Authentifizierung und Autorisierung, die Unterstützung von Draft Handling für die Zwischenspeicherung von Nutzereingaben, sowie Mechanismen für die Mandantenfähigkeit. Als Laufzeitumgebung unterstützt \gls{cap} sowohl Node.js als auch Java. In der vorliegenden Arbeit kommt die Node.js-Variante zum Einsatz, die sich durch kurze Entwicklungszyklen und eine enge Integration mit dem \gls{cds}-Werkzeugkasten auszeichnet.

Ein weiterer Vorteil von \gls{cap} ist das integrierte Entwicklungswerkzeug \texttt{cds watch}, das während der lokalen Entwicklung Dateiänderungen erkennt und den Dienst automatisch neu startet. In Verbindung mit einer SQLite-Datenbank als lokalem Ersatz für SAP HANA Cloud ermöglicht dies einen schnellen Entwicklungs- und Testzyklus ohne Cloud-Abhängigkeiten.

\subsection{SAP Fiori und Fiori Elements}
\label{subsec:fiori}
SAP Fiori bezeichnet die Designrichtlinien und das Technologierahmenwerk der SAP für moderne Benutzeroberflächen betriebswirtschaftlicher Anwendungen. Die Fiori-Designprinzipien definieren fünf zentrale Leitlinien: Anwendungen sollen \textit{rollenbasiert} (auf die Aufgaben des Nutzers zugeschnitten), \textit{adaptiv} (geräteübergreifend nutzbar), \textit{kohärent} (konsistent im Erscheinungsbild), \textit{einfach} (auf das Wesentliche fokussiert) und \textit{ansprechend} (eine positive Nutzererfahrung bietend) sein \citep{sap2024fiori}.

Fiori Elements ist ein darauf aufbauendes UI-Framework, das Benutzeroberflächen auf Grundlage von \gls{odata}-Metadaten und UI-Annotationen generiert. Der wesentliche Vorteil dieses Ansatzes besteht darin, dass keine individuelle Frontend-Entwicklung erforderlich ist. Stattdessen werden die gewünschten UI-Eigenschaften -- etwa Tabellenspalten, Filterfelder, Formularabschnitte und Navigationsstrukturen -- deklarativ über Annotationen im \gls{cds}-Modell oder in separaten Annotationsdateien definiert. Das Framework erzeugt daraus zur Laufzeit die vollständige Benutzeroberfläche.

Fiori Elements stellt verschiedene Seitentypen bereit, die als sogenannte Floorplans bezeichnet werden. Die in der Praxis am häufigsten eingesetzten Floorplans sind der \textit{List Report}, der eine tabellarische Übersicht mit Such- und Filterfunktionen anzeigt, die \textit{Object Page}, die eine detaillierte Ansicht eines einzelnen Datenobjekts mit Kopfbereich und inhaltlichen Abschnitten darstellt, sowie die \textit{Worklist}, die eine aufgabenbezogene Liste ohne erweiterte Filtermöglichkeiten bereitstellt. Diese Floorplans können miteinander kombiniert werden, um vollständige Anwendungsszenarien abzubilden.

Technisch basiert Fiori Elements auf dem Framework \gls{ui5}, dem SAP-eigenen JavaScript-Framework für die Entwicklung von Webanwendungen \citep{bince2019fiori}. \gls{ui5} implementiert das Model-View-Controller-Muster und stellt eine umfangreiche Bibliothek vorgefertigter UI-Steuerelemente bereit \citep{sap2024fiorielem}. Fiori Elements nutzt diese Steuerelemente intern und abstrahiert deren Konfiguration durch den annotationsbasierten Generierungsansatz.

\subsection{SAP Build Process Automation}
\label{subsec:sbpa}
SAP Build Process Automation (\gls{spa}) ist ein cloudbasierter Dienst auf der \gls{btp}, der die Modellierung, Automatisierung und Überwachung von Geschäftsprozessen ermöglicht \citep{freund2019bpmn}. Als Nachfolger des früheren SAP Workflow Management verfolgt \gls{spa} einen Low-Code- bzw. No-Code-Ansatz, der es auch Fachanwendern ohne tiefgreifende Programmierkenntnisse erlaubt, Prozesse zu gestalten \citep{sap2024sbpa}.

Die zentralen Konzepte von \gls{spa} umfassen Prozesse, Formulare, \gls{api}-Auslöser, Aktionen, Bedingungen und Entscheidungstabellen. Ein \textit{Prozess} bildet den übergeordneten Ablauf ab und orchestriert die einzelnen Schritte. \textit{Formulare} dienen der Darstellung von Aufgaben und Eingabemasken für menschliche Bearbeiter. Über \textit{\gls{api}-Auslöser} können Prozesse programmatisch aus externen Anwendungen heraus gestartet werden, beispielsweise über \gls{rest}-Aufrufe. \textit{Aktionen} ermöglichen den Aufruf externer Dienste und \gls{api}s innerhalb eines Prozesses. \textit{Bedingungen} steuern den Kontrollfluss, und \textit{Entscheidungstabellen} erlauben die regelbasierte Auswertung von Geschäftslogik.

Für die Bearbeitung zugewiesener Aufgaben steht den Endnutzern die Anwendung \textit{My Inbox} zur Verfügung. Über My Inbox können Genehmigungsanfragen eingesehen, genehmigt oder abgelehnt werden. Die Anwendung folgt den SAP-Fiori-Designrichtlinien und ist sowohl über den SAP Build Lobby als auch über das SAP BTP Cockpit erreichbar.

Ein wesentlicher Aspekt von \gls{spa} im Kontext dieser Arbeit ist die Möglichkeit, Prozesse über \gls{api}-Auslöser zu starten und über Rückruf-Aktionen (\textit{Callbacks}) den Status an die auslösende Anwendung zurückzumelden. Dies ermöglicht eine ereignisgesteuerte Architektur, bei der die \gls{cap}-Anwendung einen Genehmigungsprozess auslöst und anschließend über das Ergebnis der Genehmigung informiert wird.

\subsection{OData-Protokoll}
\label{subsec:odata}
Das Open Data Protocol (\gls{odata}) ist ein offener Standard der Organisation OASIS (Organization for the Advancement of Structured Information Standards) für den Aufbau und die Nutzung von \gls{rest}ful \gls{api}s. In der aktuellen Version 4.0 definiert \gls{odata} ein standardisiertes Protokoll für \gls{crud}-Operationen auf Datenressourcen über das HTTP-Protokoll \citep{oasis2014odata}.

Ein \gls{odata}-Dienst beschreibt sein Datenmodell über ein maschinenlesbares Metadaten-Dokument, das unter dem Pfad \texttt{\$metadata} abrufbar ist. Dieses Dokument enthält die Definition aller Entitätstypen, deren Eigenschaften, Beziehungen untereinander (sogenannte Navigationseigenschaften) sowie die verfügbaren Operationen. Clients können dieses Metadaten-Dokument nutzen, um die Struktur des Dienstes zur Laufzeit zu ermitteln und generische Benutzeroberflächen zu erzeugen -- ein Prinzip, das von Fiori Elements intensiv genutzt wird.

Das \gls{odata}-Protokoll folgt dabei den Prinzipien von \gls{rest}, wie sie grundlegend von \citet{tilkov2015rest} beschrieben werden. Für die Abfrage von Daten stellt \gls{odata} standardisierte Abfrageoptionen bereit. Mit \texttt{\$filter} können Ergebnismengen nach bestimmten Kriterien eingeschränkt werden, \texttt{\$expand} ermöglicht das Nachladen verknüpfter Entitäten in einer einzelnen Anfrage, \texttt{\$select} beschränkt die zurückgegebenen Eigenschaften auf eine definierte Teilmenge und \texttt{\$orderby} legt die Sortierreihenfolge fest. Diese Abfrageoptionen werden als URL-Parameter übergeben und erlauben es dem Client, präzise zu steuern, welche Daten in welchem Umfang vom Server abgerufen werden.

Im Zusammenspiel mit dem \gls{cap} ergibt sich ein besonderer Vorteil: Das Framework generiert aus den \gls{cds}-Modellen automatisch vollständige \gls{odata}-V4-Dienste, sodass die Einhaltung des Standards ohne manuellen Implementierungsaufwand gewährleistet ist. Die \gls{odata}-Schnittstelle bildet damit die zentrale Kommunikationsschicht zwischen der serverseitigen Geschäftslogik und der clientseitigen Benutzeroberfläche.

\subsection{Multi-Target Applications}
\label{subsec:mta}
Das Konzept der Multi-Target Applications (\gls{mta}) adressiert die Herausforderung, dass Cloud-Anwendungen typischerweise aus mehreren, voneinander abhängigen Modulen bestehen, die als eine logische Einheit bereitgestellt werden müssen. Eine \gls{mta} fasst verschiedene Anwendungsmodule -- etwa serverseitige Dienste, Datenbankbereitstellungen und Benutzeroberflächen -- in einem gemeinsamen Archiv zusammen und beschreibt deren Abhängigkeiten und Konfigurationen in einem zentralen Deskriptor \citep{sap2024mta}.

Der \gls{mta}-Deskriptor ist eine YAML-Datei mit dem Namen \texttt{mta.yaml}, die im Wurzelverzeichnis des Projekts liegt. In dieser Datei werden die einzelnen Module, ihre Typen (z.\,B. \texttt{nodejs} für serverseitige Anwendungen, \texttt{hdb} für \gls{hdi}-Datenbankcontainer oder \texttt{approuter.nodejs} für den Application Router), ihre Abhängigkeiten untereinander sowie die benötigten Plattformdienste deklariert. Durch die zentrale Beschreibung aller Bestandteile wird sichergestellt, dass bei der Bereitstellung sämtliche Module in der korrekten Reihenfolge und mit der richtigen Konfiguration installiert werden.

Für die Erstellung des Bereitstellungsarchivs kommt das MTA Build Tool (\texttt{mbt}) zum Einsatz, das aus dem Quellcode und dem Deskriptor ein \texttt{.mtar}-Archiv erzeugt. Dieses Archiv wird anschließend mit dem \gls{cf}-CLI-Plugin \texttt{cf deploy} auf die \gls{btp} übertragen. Der Deployer auf der Plattform interpretiert den Deskriptor, erstellt die benötigten Service-Instanzen, stellt die Module bereit und konfiguriert die Verbindungen zwischen den Komponenten.

Im Kontext der vorliegenden Arbeit besteht die \gls{mta} aus drei Modulen: dem \gls{cap}-Serverdienst, der die Geschäftslogik und die \gls{odata}-Schnittstelle bereitstellt, dem \gls{hdi}-Datenbankbereitsteller, der die Datenbankartefakte in SAP HANA Cloud installiert, sowie dem Application Router, der die Fiori-Benutzeroberfläche ausliefert und die Authentifizierung über \gls{xsuaa} koordiniert. Ergänzt werden diese Module durch die konfigurierten Plattformdienste für Datenbank, Authentifizierung und Destination-Verwaltung.

\section{Konzeption}
\label{sec:konzeption}

Aufbauend auf den im vorherigen Kapitel erarbeiteten technologischen Grundlagen wird in diesem Kapitel die Konzeption der Anwendung zur Bestellanforderungsfreigabe auf der \gls{btp} dargelegt. Zunächst werden die funktionalen und nicht-funktionalen Anforderungen systematisch erfasst und priorisiert. Darauf aufbauend wird der Architekturentwurf vorgestellt, der die Schichtenstruktur des Systems beschreibt. Anschließend werden das Datenmodell, das Prozessdesign des Genehmigungsworkflows sowie das Sicherheitskonzept erläutert.

\subsection{Anforderungsanalyse}
\label{subsec:anforderungsanalyse}

Die Anforderungsanalyse bildet die Grundlage für den Entwurf und die spätere Implementierung des Systems \citep{dumas2018bpm}. Die erhobenen Anforderungen lassen sich in funktionale und nicht-funktionale Anforderungen untergliedern. Funktionale Anforderungen beschreiben das erwartete Verhalten des Systems aus Sicht der Anwender, während nicht-funktionale Anforderungen Qualitätsmerkmale und Rahmenbedingungen der technischen Umsetzung spezifizieren.

\subsubsection{Funktionale Anforderungen}
\label{subsubsec:funktionale-anforderungen}

Die funktionalen Anforderungen wurden aus dem Anwendungsfall eines unternehmensinternen Beschaffungsprozesses abgeleitet. Zentrales Ziel ist es, den Lebenszyklus einer Bestellanforderung -- von der Erstellung über die Genehmigung bis zur Statusrückmeldung -- vollständig abzubilden. Tabelle~\ref{tab:funktionale-anforderungen} fasst die identifizierten funktionalen Anforderungen zusammen.

\begin{table}[H]
  \centering
  \caption{Funktionale Anforderungen}
  \label{tab:funktionale-anforderungen}
  \begin{tabularx}{\textwidth}{l X l}
    \hline
    \textbf{ID} & \textbf{Beschreibung} & \textbf{Priorität} \\
    \hline
    FA1 & Benutzer können Bestellanforderungen mit zugehörigen Positionen erstellen, bearbeiten und löschen (\gls{crud}-Operationen). & Hoch \\
    FA2 & Jede Bestellanforderung enthält eine Beschreibung, einen Anforderer, eine Kostenstelle sowie beliebig viele Positionen mit Materialnummer, Beschreibung, Menge und Einzelpreis. & Hoch \\
    FA3 & Nettobeträge der einzelnen Positionen sowie der Gesamtbetrag der Bestellanforderung werden automatisch berechnet. & Hoch \\
    FA4 & Benutzer können eine Bestellanforderung über die Aktion \glqq Freigeben\grqq{} zur Genehmigung einreichen. & Hoch \\
    FA5 & Die Freigabe löst einen Genehmigungsworkflow in \gls{spa} aus, der über einen \gls{api}-Trigger gestartet wird. & Hoch \\
    FA6 & Der Genehmiger kann die Bestellanforderung in der My-Inbox-Anwendung von \gls{spa} genehmigen oder ablehnen. & Hoch \\
    FA7 & Nach der Entscheidung des Genehmigers wird der Status der Bestellanforderung über einen Callback-Endpunkt in der \gls{cap}-Anwendung aktualisiert. & Hoch \\
    FA8 & Unvollständige Bestellanforderungen können als Entwurf gespeichert und zu einem späteren Zeitpunkt fortgesetzt werden (Draft-Unterstützung). & Mittel \\
    FA9 & Jede Statusänderung einer Bestellanforderung wird in einer Statushistorie protokolliert, einschließlich Zeitstempel, auslösender Quelle und optionalem Kommentar. & Mittel \\
    FA10 & Benutzer können Dokumente und Bilder (z.\,B. PDF, PNG) als Anhänge zu einer Bestellanforderung hochladen, herunterladen und löschen. Die Anhänge werden im Genehmigungsworkflow als Metadaten mit direkten Download-URLs bereitgestellt. & Mittel \\
    FA11 & Eine bereits genehmigte oder abgelehnte Bestellanforderung kann erneut bearbeitet werden. Beim Speichern wird der Status automatisch auf \texttt{Draft} zurückgesetzt, sodass eine erneute Freigabe möglich ist. & Mittel \\
    \hline
  \end{tabularx}
\end{table}

Die Anforderungen FA1 bis FA7 wurden mit der Priorität \glqq Hoch\grqq{} eingestuft, da sie den Kernprozess der Bestellanforderungsfreigabe abbilden und für eine funktionsfähige Anwendung unverzichtbar sind. Die Draft-Unterstützung (FA8), die Statushistorie (FA9), die Dokumentenanhänge (FA10) sowie die Wiederbearbeitbarkeit (FA11) erhielten die Priorität \glqq Mittel\grqq{}, da sie die Benutzerfreundlichkeit und Nachvollziehbarkeit signifikant erhöhen, jedoch nicht für den grundlegenden Genehmigungsprozess erforderlich sind.

\subsubsection{Nicht-funktionale Anforderungen}
\label{subsubsec:nichtfunktionale-anforderungen}

Neben den funktionalen Anforderungen wurden nicht-funktionale Anforderungen definiert, die die technischen Rahmenbedingungen und Qualitätsmerkmale des Systems festlegen. Diese sind in Tabelle~\ref{tab:nichtfunktionale-anforderungen} aufgeführt.

\begin{table}[H]
  \centering
  \caption{Nicht-funktionale Anforderungen}
  \label{tab:nichtfunktionale-anforderungen}
  \begin{tabularx}{\textwidth}{l X l}
    \hline
    \textbf{ID} & \textbf{Beschreibung} & \textbf{Priorität} \\
    \hline
    NFA1 & Die Anwendung wird cloud-nativ auf der \gls{btp} bereitgestellt und nutzt die \gls{cf}-Laufzeitumgebung. & Hoch \\
    NFA2 & Die Authentifizierung und Autorisierung erfolgt über den \gls{xsuaa}-Service der \gls{btp}. & Hoch \\
    NFA3 & Die Benutzeroberfläche ist responsiv und unterstützt die Nutzung auf Desktop-, Tablet- und Mobilgeräten. & Mittel \\
    NFA4 & Die Kommunikation zwischen den Komponenten erfolgt über \gls{rest}ful \gls{api}s und standardisierte Protokolle (\gls{odata} V4). & Hoch \\
    \hline
  \end{tabularx}
\end{table}

\subsection{Architekturentwurf}
\label{subsec:architekturentwurf}

Der Architekturentwurf orientiert sich an einer klassischen Dreischichtarchitektur, die in eine Datenbank-, eine Service- und eine Präsentationsschicht untergliedert ist. Diese Schichtung ermöglicht eine klare Trennung der Verantwortlichkeiten und erleichtert sowohl die Wartbarkeit als auch die Erweiterbarkeit des Systems. Abbildung~\ref{fig:architektur} zeigt die Gesamtarchitektur der Lösung einschließlich der Integration mit \gls{spa}.

\begin{figure}[H]
  \centering
  \begin{tikzpicture}[
    box/.style={rectangle, draw, minimum width=2.8cm, minimum height=1cm, align=center, font=\small},
    layer/.style={rectangle, draw, dashed, inner sep=8pt},
    arrow/.style={-{Stealth}, thick},
    node distance=0.8cm
  ]
    % CAP Application Layers
    \node[box, fill=blue!15] (ui) {Fiori Elements\\(List Report, Object Page)};
    \node[box, fill=green!15, below=of ui] (srv) {CAP Service\\(Node.js, OData V4)};
    \node[box, fill=orange!15, below=of srv] (db) {SAP HANA Cloud\\(HDI Container)};

    % Layer labels
    \node[left=0.3cm of ui, font=\scriptsize\itshape] {Präsentation};
    \node[left=0.3cm of srv, font=\scriptsize\itshape] {Service};
    \node[left=0.3cm of db, font=\scriptsize\itshape] {Datenbank};

    % SPA
    \node[box, fill=purple!15, right=2.5cm of srv] (spa) {SAP Build\\Process Automation};
    \node[box, fill=purple!10, below=0.5cm of spa] (inbox) {My Inbox};

    % Arrows within CAP
    \draw[arrow, <->] (ui) -- (srv);
    \draw[arrow, <->] (srv) -- (db);

    % Arrows to SPA
    \draw[arrow] (srv.east) -- node[above, font=\scriptsize] {API-Trigger} (spa.west);
    \draw[arrow] (spa.west) -- +(-0.5,0) |- node[below, font=\scriptsize, pos=0.25] {Callback} ([yshift=-3mm]srv.east);

    % Background boxes
    \begin{scope}[on background layer]
      \node[layer, fill=gray!5, fit=(ui)(srv)(db), label={[font=\small\bfseries]above:CAP-Anwendung}] {};
      \node[layer, fill=purple!5, fit=(spa)(inbox), label={[font=\small\bfseries]above:Workflow}] {};
    \end{scope}
  \end{tikzpicture}
  \caption{Architektur der Bestellanforderungsanwendung mit Workflow-Integration}
  \label{fig:architektur}
\end{figure}

Im Folgenden werden die einzelnen Schichten und ihre Verantwortlichkeiten erläutert.

\subsubsection{Datenbankschicht}

Die Datenbankschicht bildet das Fundament der Anwendung und ist für die persistente Speicherung aller Geschäftsdaten verantwortlich. Als Datenbanktechnologie kommt SAP HANA Cloud zum Einsatz \citep{sap2024hana}, die über einen \gls{hdi}-Container bereitgestellt wird. Das Datenbankschema wird deklarativ mithilfe von \gls{cds} definiert und bei der Bereitstellung automatisch in die entsprechenden HANA-Artefakte transformiert. Dieser Ansatz stellt sicher, dass das Datenmodell versioniert und reproduzierbar in verschiedenen Umgebungen bereitgestellt werden kann.

\subsubsection{Service-Schicht}

Die Service-Schicht stellt den zentralen Verarbeitungskern der Anwendung dar. Sie wird als \gls{cap}-Anwendung auf Basis von Node.js implementiert und exponiert die Geschäftslogik als \gls{odata}-V4-Service. Die wesentlichen Aufgaben der Service-Schicht umfassen die Bereitstellung von \gls{crud}-Operationen für Bestellanforderungen und deren Positionen, die automatische Berechnung von Nettobeträgen und Gesamtbeträgen, die Validierung der Eingabedaten bei der Freigabe einer Bestellanforderung, den Start des Genehmigungsworkflows in \gls{spa} über die \gls{api}-Schnittstelle sowie die Verarbeitung des Callback-Aufrufs nach Abschluss des Genehmigungsprozesses. Die Anbindung an \gls{spa} erfolgt über einen dedizierten Workflow-Client, der den OAuth2 Client Credentials Flow zur Authentifizierung nutzt und die Kommunikation mit der \gls{spa}-\gls{api} kapselt.

\subsubsection{Präsentationsschicht}

Die Präsentationsschicht besteht aus zwei Komponenten. Die primäre Benutzeroberfläche wird mithilfe von SAP Fiori Elements realisiert und basiert auf dem List-Report- und Object-Page-Floorplan. Über deklarative \gls{cds}-Annotationen wird die Benutzeroberfläche direkt aus dem Service-Modell generiert, wodurch ein konsistentes und SAP-konformes Benutzererlebnis gewährleistet wird. Die zweite Komponente der Präsentationsschicht ist die My-Inbox-Anwendung von \gls{spa}, in der Genehmiger ihre Aufgaben einsehen und bearbeiten können.

\subsection{Datenmodell}
\label{subsec:datenmodell}

Das Datenmodell bildet die zentrale Informationsstruktur der Anwendung ab und wird deklarativ in \gls{cds} definiert. Es besteht aus zwei Entitäten, die über eine Kompositionsbeziehung miteinander verknüpft sind: \texttt{PurchaseRequisitions} als übergeordnete Entität und \texttt{Items} als untergeordnete Positionsentität.

Tabelle~\ref{tab:entity-purchaserequisitions} beschreibt die Attribute der Entität \texttt{PurchaseRequisitions}, die eine einzelne Bestellanforderung repräsentiert.

\begin{table}[H]
  \centering
  \caption{Attribute der Entität PurchaseRequisitions}
  \label{tab:entity-purchaserequisitions}
  \begin{tabularx}{\textwidth}{l l X l}
    \hline
    \textbf{Attribut} & \textbf{Datentyp} & \textbf{Beschreibung} & \textbf{Constraint} \\
    \hline
    ID            & UUID            & Eindeutiger Bezeichner       & Primärschlüssel \\
    description   & String(256)     & Bezeichnung der Anforderung  & Pflichtfeld \\
    requester     & String(128)     & Name des Anforderers         & Pflichtfeld \\
    costCenter    & String(10)      & Zugeordnete Kostenstelle     & Pflichtfeld \\
    totalAmount   & Decimal(15,2)   & Berechneter Gesamtbetrag     & Berechnet \\
    currency      & Currency        & Währungsschlüssel            & -- \\
    status        & String(20)      & Prozessstatus                & Standard: Draft \\
    workflowId    & String(256)     & ID der Workflow-Instanz      & -- \\
    approverComment & String(1000)  & Kommentar des Genehmigers    & -- \\
    approvedBy    & String(128)     & Name des Entscheiders        & -- \\
    decidedAt     & Timestamp       & Zeitpunkt der Entscheidung   & -- \\
    items         & Composition     & Zugehörige Positionen        & 1:n \\
    statusHistory & Composition     & Protokollierte Statusänderungen & 1:n \\
    attachments   & Composition     & Hochgeladene Dokumentenanhänge & 1:n \\
    \hline
  \end{tabularx}
\end{table}

Das Attribut \texttt{status} kann die Werte \texttt{Draft}, \texttt{Submitted}, \texttt{Approved} und \texttt{Rejected} annehmen und bildet somit den vollständigen Lebenszyklus einer Bestellanforderung ab. Das Feld \texttt{workflowId} dient der Zuordnung einer Bestellanforderung zu ihrer zugehörigen Workflow-Instanz in \gls{spa} und wird erst beim Einreichen der Anforderung befüllt. Die Felder \texttt{approverComment}, \texttt{approvedBy} und \texttt{decidedAt} speichern die Genehmigungsinformationen, die über den Callback von \gls{spa} zurückgemeldet werden.

Tabelle~\ref{tab:entity-items} beschreibt die Attribute der Entität \texttt{Items}, die eine einzelne Position innerhalb einer Bestellanforderung darstellt.

\begin{table}[H]
  \centering
  \caption{Attribute der Entität Items}
  \label{tab:entity-items}
  \begin{tabularx}{\textwidth}{l l X l}
    \hline
    \textbf{Attribut} & \textbf{Datentyp} & \textbf{Beschreibung} & \textbf{Constraint} \\
    \hline
    ID            & UUID            & Eindeutiger Bezeichner       & Primärschlüssel \\
    parent        & Association     & Verweis auf die übergeordnete Bestellanforderung & Fremdschlüssel \\
    materialNo    & String(40)      & Materialnummer               & Pflichtfeld \\
    description   & String(256)     & Beschreibung der Position    & -- \\
    quantity      & Decimal(13,3)   & Bestellmenge                 & Pflichtfeld \\
    unitPrice     & Decimal(15,2)   & Einzelpreis pro Mengeneinheit & Pflichtfeld \\
    currency      & Currency        & Währungsschlüssel            & -- \\
    netAmount     & Decimal(15,2)   & Berechneter Nettobetrag      & Berechnet \\
    \hline
  \end{tabularx}
\end{table}

Die Beziehung zwischen \texttt{PurchaseRequisitions} und \texttt{Items} ist als Komposition modelliert. Dies bedeutet, dass Positionen existenzabhängig von ihrer übergeordneten Bestellanforderung sind: Wird eine Bestellanforderung gelöscht, werden automatisch auch alle zugehörigen Positionen entfernt. Der Nettobetrag einer Position ergibt sich aus der Multiplikation von Menge (\texttt{quantity}) und Einzelpreis (\texttt{unitPrice}). Der Gesamtbetrag der Bestellanforderung wird als Summe aller Positionsnettobeträge berechnet.

Neben den Positionen verfügt die Bestellanforderung über zwei weitere Kompositionen. Die Entität \texttt{StatusHistory} protokolliert jeden Statuswechsel mit dem alten und neuen Status, einem Zeitstempel, dem auslösenden Benutzer, einem optionalen Kommentar sowie der Quelle der Änderung (\texttt{USER}, \texttt{WORKFLOW} oder \texttt{CALLBACK}). Diese Protokollierung ermöglicht eine lückenlose Nachvollziehbarkeit des Genehmigungsprozesses.

Die Komposition \texttt{attachments} wird über das offizielle CAP-Plugin \texttt{@cap-js/attachments} realisiert. Dieses Plugin stellt eine vorkonfigurierte \texttt{Attachments}-Entität bereit, die Dateiname, MIME-Typ und den binären Inhalt als BLOB in der HANA-Datenbank speichert. Die Konfiguration als datenbankbasierter Speicher (\texttt{kind: db}) macht einen externen Object Store überflüssig. Durch die Kompositionsbeziehung werden die Anhänge im Draft-Modus gemeinsam mit der übergeordneten Bestellanforderung verwaltet und bei Fiori Elements automatisch als Upload-Bereich in der Object Page dargestellt.

\subsection{Prozessdesign}
\label{subsec:prozessdesign}

Der Genehmigungsprozess für Bestellanforderungen erstreckt sich über zwei Systeme -- die \gls{cap}-Anwendung und \gls{spa} -- und durchläuft mehrere definierte Phasen. Im Folgenden wird der End-to-End-Prozess beschrieben.

\subsubsection{Ablauf des Genehmigungsprozesses}

Der Prozess beginnt mit der Erstellung einer Bestellanforderung durch den Anforderer in der Fiori-Elements-Oberfläche. Die Anforderung wird zunächst im Status \texttt{Draft} angelegt und kann beliebig bearbeitet werden. Der Anforderer fügt Positionen mit Materialnummer, Menge und Einzelpreis hinzu, wobei die Nettobeträge und der Gesamtbetrag automatisch berechnet werden.

Sobald die Bestellanforderung vollständig erfasst ist, löst der Anforderer über die Aktion \glqq Freigeben\grqq{} den Genehmigungsprozess aus. Dabei werden zunächst serverseitige Validierungen durchgeführt: Es wird geprüft, ob sich die Anforderung nicht bereits im Status \texttt{Submitted} befindet, ob mindestens eine Position vorhanden ist und ob der Gesamtbetrag größer als null ist. Nach erfolgreicher Validierung wird der Status auf \texttt{Submitted} gesetzt und ein Genehmigungsworkflow in \gls{spa} über die Workflow-\gls{api} gestartet. Die zurückgegebene Workflow-Instanz-ID wird in der Bestellanforderung gespeichert.

Der Genehmiger erhält die Aufgabe in der My-Inbox-Anwendung von \gls{spa}. Dort werden die relevanten Informationen der Bestellanforderung -- Beschreibung, Anforderer und Gesamtbetrag -- angezeigt. Der Genehmiger kann die Anforderung genehmigen oder ablehnen.

Nach der Entscheidung des Genehmigers ruft \gls{spa} den Callback-Endpunkt der \gls{cap}-Anwendung auf und übermittelt die Workflow-ID, den Entscheidungsstatus, den Namen des Genehmigers sowie einen optionalen Kommentar. Die \gls{cap}-Anwendung identifiziert die zugehörige Bestellanforderung anhand der Workflow-ID und aktualisiert den Status auf \texttt{Approved} oder \texttt{Rejected}. Die Genehmigungsinformationen werden in der Entität gespeichert und in der Benutzeroberfläche angezeigt. Jede Statusänderung wird in der Statushistorie protokolliert. Sollte der Workflow-Start fehlschlagen, wird der Status automatisch auf \texttt{Draft} zurückgesetzt, sodass der Anforderer den Vorgang erneut auslösen kann.

Darüber hinaus unterstützt der Prozess die erneute Bearbeitung nach einer Genehmigung oder Ablehnung. Wird eine Bestellanforderung im Status \texttt{Approved} oder \texttt{Rejected} bearbeitet und gespeichert, setzt die Anwendung den Status automatisch auf \texttt{Draft} zurück und entfernt die vorherigen Genehmigungsinformationen. Auf diese Weise kann der Anforderer die Anforderung überarbeiten und erneut zur Genehmigung einreichen.

\subsubsection{Workflow-Design in SAP Build Process Automation}

Abbildung~\ref{fig:workflow} zeigt den in \gls{spa} modellierten Genehmigungsprozess als Flussdiagramm.

\begin{figure}[H]
  \centering
  \begin{tikzpicture}[
    startstop/.style={circle, draw, fill=green!20, minimum size=0.8cm},
    process/.style={rectangle, draw, fill=blue!15, minimum width=3cm, minimum height=0.8cm, align=center, font=\small},
    decision/.style={diamond, draw, fill=yellow!20, minimum width=2cm, minimum height=1.5cm, align=center, font=\small, aspect=2},
    arrow/.style={-{Stealth}, thick},
    node distance=1.2cm
  ]
    % Start
    \node[startstop] (start) {};
    \node[left=0.2cm of start, font=\scriptsize] {Start};

    % API Trigger
    \node[process, below=of start] (trigger) {API-Trigger\\(Workflow starten)};

    % User Task
    \node[process, below=of trigger] (task) {Benutzeraufgabe\\(Genehmigungsformular)};

    % Decision
    \node[decision, below=1.5cm of task] (decision) {Entscheidung?};

    % Callback nodes
    \node[process, below left=1.2cm and 0.8cm of decision] (approved) {Callback:\\Status = Approved};
    \node[process, below right=1.2cm and 0.8cm of decision] (rejected) {Callback:\\Status = Rejected};

    % End
    \node[startstop, below=2.5cm of decision, fill=red!20] (end) {};
    \node[left=0.2cm of end, font=\scriptsize] {Ende};

    % Arrows
    \draw[arrow] (start) -- (trigger);
    \draw[arrow] (trigger) -- (task);
    \draw[arrow] (task) -- (decision);
    \draw[arrow] (decision) -| node[above, font=\scriptsize, pos=0.25] {Genehmigen} (approved);
    \draw[arrow] (decision) -| node[above, font=\scriptsize, pos=0.25] {Ablehnen} (rejected);
    \draw[arrow] (approved) |- (end);
    \draw[arrow] (rejected) |- (end);
  \end{tikzpicture}
  \caption{Genehmigungsworkflow in SAP Build Process Automation}
  \label{fig:workflow}
\end{figure}

Der Prozess folgt einer linearen Struktur. Zunächst wird der Prozess durch einen \gls{api}-Aufruf aus der \gls{cap}-Anwendung gestartet, wobei der Kontextparameter die Bestellanforderungs-ID, die Beschreibung, den Anforderer, die Kostenstelle, den Gesamtbetrag, die einzelnen Positionen, die Metadaten der Anhänge sowie einen Deep Link zur Fiori-Anwendung enthält. Anschließend wird eine Benutzeraufgabe erstellt und dem Genehmiger in der My Inbox zugewiesen, wo das Formular die Kontextdaten anzeigt und die Optionen \glqq Genehmigen\grqq{} und \glqq Ablehnen\grqq{} bietet. Basierend auf der Entscheidung wird der Prozessfluss verzweigt. Unabhängig von der Entscheidung wird ein HTTP-Aufruf an den Callback-Endpunkt der \gls{cap}-Anwendung ausgeführt, der die Workflow-ID, den Entscheidungsstatus, den Namen des Genehmigers und einen optionalen Kommentar übermittelt. Abschließend wird der Workflow beendet.

\subsection{Sicherheitskonzept}
\label{subsec:sicherheitskonzept}

Das Sicherheitskonzept der Anwendung adressiert die Authentifizierung, die Autorisierung sowie die sichere Kommunikation zwischen den beteiligten Diensten auf der \gls{btp}.

\subsubsection{Authentifizierung und Autorisierung}

Die Authentifizierung der Benutzer erfolgt über den \gls{xsuaa}-Service der \gls{btp} \citep{sap2024xsuaa}. Dieser fungiert als OAuth2-Autorisierungsserver und stellt \gls{jwt}-basierte Zugriffstoken aus, die bei jedem Zugriff auf die \gls{cap}-Anwendung validiert werden. Die Autorisierung basiert auf einem rollenbasierten Zugriffsmodell mit zwei Rollen: Die Rolle \textit{Requester} berechtigt zum Erstellen, Bearbeiten und Einreichen von Bestellanforderungen, während die Rolle \textit{Approver} zum Genehmigen oder Ablehnen von Bestellanforderungen in der My Inbox von \gls{spa} berechtigt.

\subsubsection{Service-zu-Service-Kommunikation}

Die Kommunikation zwischen der \gls{cap}-Anwendung und \gls{spa} erfolgt über den OAuth2 Client Credentials Flow. Dabei authentifiziert sich die \gls{cap}-Anwendung mit einer Client-ID und einem Client-Secret beim Token-Endpunkt des \gls{xsuaa}-Service und erhält ein Zugriffstoken, das für die nachfolgenden \gls{api}-Aufrufe verwendet wird. Dieser Mechanismus gewährleistet, dass ausschließlich autorisierte Dienste den Genehmigungsworkflow starten können.

Für die Anbindung externer Dienste werden BTP-Destinations genutzt. Diese kapseln die Verbindungsinformationen -- einschließlich der Endpunkt-URL und der Authentifizierungsparameter -- und ermöglichen eine zentrale Verwaltung der Service-Konnektivität. Durch die Trennung von Anwendungslogik und Konfigurationsdaten wird die Sicherheit erhöht, da sensible Zugangsdaten nicht im Quellcode hinterlegt werden müssen.

\section{Implementierung}

Dieses Kapitel beschreibt die technische Umsetzung des Bestellanforderungsprozesses. Die Implementierung erfolgt auf Basis des \gls{cap} mit Node.js, SAP Fiori Elements und \gls{spa}. Vollständige Code-Listings befinden sich im Anhang.

\subsection{Projektstruktur}

Die Anwendung folgt der \gls{cap}-Konvention mit drei Hauptverzeichnissen: \texttt{db/} enthält das \gls{cds}-Domänenmodell, \texttt{srv/} die Servicedefinition und Geschäftslogik, \texttt{app/} die Fiori-Elements-Anwendung. Die Dateien \texttt{mta.yaml} und \texttt{xs-security.json} steuern Deployment und Sicherheit.

\subsection{Datenbankschicht}

Das Domänenmodell definiert zwei Hauptentitäten in \gls{cds} (siehe Anhang~A). Die Entität \texttt{PurchaseRequisitions} repräsentiert Bestellanforderungen mit Beschreibung, Anforderer, Kostenstelle, Gesamtbetrag und Status (\texttt{Draft}, \texttt{Submitted}, \texttt{Approved}, \texttt{Rejected}). Die Entität \texttt{Items} enthält die zugehörigen Positionen mit Materialnummer, Menge, Einzelpreis und berechnetem Nettobetrag. Die Beziehung ist als Komposition modelliert, sodass Positionen gemeinsam mit der übergeordneten Anforderung verwaltet werden. Zusätzlich protokolliert eine \texttt{StatusHistory}-Entität alle Statusänderungen, und das Plugin \texttt{@cap-js/attachments} ermöglicht Dokumentenanhänge.

\subsection{Service-Schicht}

Die Servicedefinition (siehe Anhang~B) exponiert das Domänenmodell als \gls{odata}-V4-Service mit drei Kernfunktionen: Die Draft-Unterstützung durch die Annotation \texttt{@odata.draft.enabled} ermöglicht das schrittweise Bearbeiten von Anforderungen. Die Submit-Aktion validiert die Anforderung und startet den Genehmigungsworkflow. Der Callback-Endpunkt empfängt die Genehmigungsentscheidung von \gls{spa}. Die Geschäftslogik in den Event-Handlern berechnet automatisch Netto- und Gesamtbeträge beim Speichern und protokolliert alle Statusänderungen in der Historie.

\subsection{Workflow-Integration}

Die Kommunikation mit \gls{spa} erfolgt über einen dedizierten \texttt{WorkflowClient}, der den OAuth2 Client Credentials Flow implementiert. Beim Freigeben einer Anforderung wird ein Genehmigungsworkflow gestartet, der folgende Kontextdaten erhält: Anforderungs-ID, Beschreibung, Anforderer, Gesamtbetrag, Positionsliste, Anhänge-URLs sowie ein Deep Link zur Fiori-Anwendung.

Nach der Genehmigungsentscheidung ruft \gls{spa} den Callback-Endpunkt auf und übermittelt Status, Genehmigername und Kommentar. Die CAP-Anwendung aktualisiert daraufhin den Status der Bestellanforderung.

\subsection{Frontend}

Die Benutzeroberfläche basiert vollständig auf SAP Fiori Elements und wird deklarativ über \gls{cds}-Annotationen konfiguriert. Der List Report zeigt alle Bestellanforderungen tabellarisch an, die Object Page ermöglicht die Detailbearbeitung mit Sektionen für Kopfdaten, Positionen, Genehmigungsinformationen, Statusverlauf und Anhänge.

Eine benutzerdefinierte Schaltfläche ermöglicht die Navigation zur SPA My Inbox, wo Genehmiger ihre Aufgaben bearbeiten können.

\subsection{Deployment}

Die Bereitstellung erfolgt als \gls{mta} auf Cloud Foundry (siehe Anhang~D). Das Deployment umfasst drei Module: Der CAP-Server ist eine Node.js-Anwendung mit \gls{odata}-Service und Geschäftslogik. Der HDI-Deployer überführt das CDS-Modell in das HANA-Datenbankschema. Der Application Router liefert die Fiori-Oberfläche aus und koordiniert die Authentifizierung. Die \gls{spa}-Konfiguration erfolgt über Umgebungsvariablen, was die Trennung von Code und Konfiguration nach dem Twelve-Factor-App-Prinzip gewährleistet \citep{wiggins2017twelvefactor}.

\section{Ergebnisse und Evaluation}

Dieses Kapitel stellt die Ergebnisse der prototypischen Implementierung vor und bewertet diese anhand der definierten Anforderungen.

\subsection{Funktionstest}

Zur Validierung wurde ein manueller Funktionstest durchgeführt, der den vollständigen Prozess von der Anlage einer Bestellanforderung bis zum Abschluss des Genehmigungsworkflows abdeckt.

Der \textbf{Genehmigungspfad} umfasst: Anlegen einer Bestellanforderung im Status \texttt{Draft}, Erfassen von Kopfdaten und Positionen, automatische Betragsberechnung beim Speichern, Freigabe über die \textit{Freigeben}-Aktion, Start des Workflows in \gls{spa}, Genehmigung in der My Inbox und Statusaktualisierung auf \texttt{Approved} über den Callback.

Der \textbf{Ablehnungspfad} verläuft analog, wobei der Status auf \texttt{Rejected} gesetzt wird und eine erneute Bearbeitung und Einreichung möglich ist.

Tabelle~\ref{tab:testfaelle} fasst die Testfälle zusammen.

\begin{table}[H]
\centering
\caption{Durchgeführte Testfälle}
\label{tab:testfaelle}
\begin{tabularx}{\textwidth}{c|X|c}
\hline
\textbf{ID} & \textbf{Testfall} & \textbf{Ergebnis} \\
\hline
T1 & Bestellanforderung anlegen und speichern & Bestanden \\
T2 & Automatische Betragsberechnung & Bestanden \\
T3 & Freigeben ohne Positionen (Validierung) & Bestanden \\
T4 & Freigeben mit Positionen, Workflow-Start & Bestanden \\
T5 & Genehmigung in SPA My Inbox & Bestanden \\
T6 & Ablehnung in SPA My Inbox & Bestanden \\
T7 & Draft-Unterstützung & Bestanden \\
T8 & Dokumentenanhänge & Bestanden \\
T9 & Statushistorie & Bestanden \\
T10 & Wiederbearbeitung nach Entscheidung & Bestanden \\
\hline
\end{tabularx}
\end{table}

\subsection{Anforderungsabgleich}

Alle elf funktionalen Anforderungen aus Kapitel~3 wurden vollständig umgesetzt: \gls{crud}-Operationen (FA1), Positionserfassung mit Betragsberechnung (FA2, FA3), Statusverwaltung (FA4), Freigabeprozess mit Validierung (FA5), Workflow-Integration (FA6, FA7), Draft-Unterstützung (FA8), Statushistorie (FA9), Dokumentenanhänge (FA10) und Wiederbearbeitbarkeit (FA11).

\subsection{Bewertung}

\paragraph{Stärken}
Die Kombination von \gls{cap} und Fiori Elements ermöglicht eine hohe Entwicklungsgeschwindigkeit durch deklarative Modellierung. Die Benutzeroberfläche wird vollständig annotationsgesteuert generiert, ohne manuellen Frontend-Code. Die durchgängige Integration innerhalb der \gls{btp} und die klare Schichtentrennung fördern Wartbarkeit und Erweiterbarkeit.

\paragraph{Herausforderungen}
Technische Herausforderungen umfassten die Feldnamen-Konvertierung zwischen \gls{cap} und \gls{spa}, die regionenübergreifende Destination-Konfiguration sowie die Betragsberechnung im Zusammenspiel mit dem Draft-Mechanismus.

\paragraph{Limitierungen}
Der Prototyp ist als Einzelmandantenlösung konzipiert, verfügt über keine mehrstufigen Genehmigungsregeln und keine automatisierten Tests. Für einen produktiven Einsatz wären diese Erweiterungen erforderlich.

\section{Fazit und Ausblick}
\label{sec:fazit}

Das abschließende Kapitel fasst die wesentlichen Ergebnisse der vorliegenden Arbeit zusammen, beantwortet die eingangs formulierte Fragestellung und zeigt Ansatzpunkte für eine Weiterentwicklung des vorgestellten Prototyps auf.

\subsection{Zusammenfassung}

Gegenstand dieser Arbeit war die Konzeption und prototypische Implementierung eines Bestellanforderungsprozesses mit integriertem Genehmigungsworkflow auf der \gls{btp}. Ausgehend von der Problemstellung, dass Beschaffungsprozesse in vielen Unternehmen noch durch manuelle, papierbasierte oder E-Mail-gestützte Abläufe gekennzeichnet sind, wurde untersucht, wie sich durch die Kombination von \gls{cap}, SAP Fiori Elements und \gls{spa} ein durchgängiger, cloudnativer Geschäftsprozess auf der \gls{btp} realisieren lässt.

Die in der Einleitung definierten vier Teilziele wurden im Rahmen der Arbeit vollständig adressiert. Zunächst wurden in einer Anforderungsanalyse elf funktionale und vier nicht-funktionale Anforderungen systematisch erhoben und priorisiert (Teilziel~1). Darauf aufbauend wurde eine Anwendungsarchitektur entworfen, die auf einer klaren Drei-Schichten-Architektur basiert (Teilziel~2). Die Datenhaltungsschicht wird durch ein \gls{cds}-basiertes Datenmodell realisiert, das Bestellanforderungen mit ihren zugehörigen Positionen und Statusinformationen abbildet. Die Geschäftslogikschicht wurde mithilfe des \gls{cap} in Node.js implementiert und stellt die Fachlogik über \gls{odata}-Dienste bereit. Die Präsentationsschicht wurde auf Basis von SAP Fiori Elements umgesetzt, wobei der annotationsgetriebene Ansatz des \gls{cap} eine effiziente Generierung der Benutzeroberfläche ermöglichte, ohne dass umfangreicher Frontend-Code geschrieben werden musste.

Im Rahmen der prototypischen Implementierung (Teilziel~3) wurde der gesamte End-to-End-Prozess realisiert: Nach dem Erstellen einer Bestellanforderung über die Fiori-Elements-Oberfläche -- einschließlich des Hochladens von Dokumentenanhängen -- wird der Genehmigungsprozess in \gls{spa} automatisiert ausgelöst. Dem Genehmiger werden im Genehmigungsformular die vollständigen Kontextdaten einschließlich der Positionsliste, der Anhänge als Download-Links und einem Deep Link zur Fiori-Anwendung bereitgestellt. Das Ergebnis der Genehmigungsentscheidung wird über einen \gls{api}-Callback an den \gls{cap}-Dienst zurückgeschrieben, sodass der Status der Bestellanforderung in Echtzeit aktualisiert wird. Alle Statusänderungen werden in einer Statushistorie protokolliert, die den vollständigen Lebenszyklus jeder Bestellanforderung dokumentiert. Darüber hinaus können bereits genehmigte oder abgelehnte Anforderungen erneut bearbeitet und eingereicht werden.

Die Anwendung wurde als \gls{mta} auf Cloud Foundry bereitgestellt und anhand von elf Testfällen evaluiert (Teilziel~4). Dabei konnten alle elf funktionalen Anforderungen vollständig nachgewiesen werden.

Die Implementierung hat gezeigt, dass der annotationsgetriebene Ansatz des \gls{cap} in Kombination mit SAP Fiori Elements eine erhebliche Beschleunigung der Oberflächenentwicklung ermöglicht. Durch die deklarative Definition von Annotationen im \gls{cds}-Modell konnten Listen- und Detailansichten sowie Formulare ohne manuellen \gls{ui5}-Code generiert werden. Darüber hinaus hat sich \gls{spa} als geeignetes Werkzeug für die Low-Code-Modellierung von Genehmigungsworkflows erwiesen. Die grafische Prozessmodellierung ermöglicht eine schnelle Umsetzung von Workflow-Szenarien, während die nahtlose \gls{api}-Integration den bidirektionalen Datenaustausch mit der \gls{cap}-Anwendung sicherstellt.

Zusammenfassend lässt sich die eingangs formulierte Forschungsfrage positiv beantworten: Es ist sowohl technisch machbar als auch effizient, durch die Kombination von \gls{cap}, SAP Fiori Elements und \gls{spa} einen durchgängigen, cloudnativen Bestellanforderungsprozess mit Genehmigungsworkflow auf der \gls{btp} zu realisieren. Die Plattform stellt die notwendigen Bausteine bereit, um cloudnative Unternehmensanwendungen mit durchgängiger Prozessautomatisierung zu entwickeln. Die klare Trennung der Verantwortlichkeiten zwischen den einzelnen Diensten fördert dabei eine wartbare und erweiterbare Architektur, die den Anforderungen moderner Geschäftsprozessdigitalisierung gerecht wird.

\subsection{Ausblick}

Der im Rahmen dieser Arbeit entwickelte Prototyp bildet eine solide Grundlage, auf der weiterführende Entwicklungen aufbauen können. Im Folgenden werden ausgewählte Erweiterungsmöglichkeiten skizziert, die den praktischen Nutzen und die Produktionsreife der Anwendung erhöhen würden.

Eine naheliegende Erweiterung besteht in der Integration mit einem SAP S/4HANA-System \citep{metzger2021sap}. Im aktuellen Prototyp endet der Prozess mit der Genehmigung der Bestellanforderung. In einem produktiven Szenario sollte nach erfolgter Genehmigung automatisiert eine Bestellung im S/4HANA-Backend erzeugt werden, um den Beschaffungsprozess vollständig abzubilden. Hierzu könnten die von SAP bereitgestellten \gls{api}s des SAP S/4HANA-Systems genutzt werden, die sich über den SAP Destination Service in die \gls{btp}-Anwendung einbinden lassen.

Ein weiterer wesentlicher Aspekt ist die Einführung eines rollenbasierten Berechtigungskonzepts. Im vorliegenden Prototyp wird nicht zwischen der Rolle des Anfordernden und der des Genehmigenden unterschieden. In einer produktiven Umgebung sollte über den \gls{xsuaa}-Dienst eine feingranulare Rollenzuweisung erfolgen, sodass Anfordernde ausschließlich eigene Bestellanforderungen erstellen und einsehen können, während Genehmigende nur auf die ihnen zugewiesenen Genehmigungsaufgaben Zugriff erhalten.

Darüber hinaus wäre die Implementierung eines mehrstufigen Genehmigungsprozesses in Abhängigkeit vom Bestellwert von praktischem Nutzen. So könnte beispielsweise festgelegt werden, dass Bestellanforderungen ab einem bestimmten Schwellenwert -- etwa ab 10.000~EUR -- eine zusätzliche Genehmigung durch eine übergeordnete Instanz erfordern. \gls{spa} bietet mit seinen Entscheidungsregeln und Verzweigungslogiken die notwendigen Modellierungsmöglichkeiten, um derartige Szenarien umzusetzen.

Zur Verbesserung der Benutzererfahrung wäre zudem die Integration von Benachrichtigungsfunktionen sinnvoll. E-Mail- oder Push-Benachrichtigungen bei ausstehenden Genehmigungsaufgaben würden sicherstellen, dass Genehmigende zeitnah auf offene Anforderungen aufmerksam gemacht werden und Verzögerungen im Prozessablauf reduziert werden. Die im Prototyp implementierte Statushistorie bildet bereits eine Grundlage für die Nachvollziehbarkeit von Genehmigungsentscheidungen. Für ein umfassendes Compliance-Reporting wäre jedoch eine Erweiterung um die Protokollierung sämtlicher Datenänderungen an Kopfdaten und Positionen empfehlenswert.

Hinsichtlich der Bereitstellung und des Benutzerzugangs bietet sich eine Migration zu SAP Build Work Zone an. Durch die Einbindung der Anwendung in ein zentrales Launchpad könnten Nutzende über eine einheitliche Oberfläche auf die Bestellanforderungsanwendung sowie auf weitere Geschäftsanwendungen zugreifen, was die Akzeptanz und die Integration in bestehende Arbeitsabläufe verbessern würde.

Schließlich sollte für eine produktive Nutzung die Qualitätssicherung durch automatisierte Tests gestärkt werden. Das \gls{cap} stellt hierfür ein eigenes Testframework bereit \citep{herger2020cap}, mit dem sowohl Unit-Tests für einzelne Geschäftslogikkomponenten als auch Integrationstests für die \gls{odata}-Schnittstellen umgesetzt werden können. Eine Optimierung der Anwendung für große Datenvolumina, etwa durch serverseitige Paginierung und gezielte Datenbankindizierung, wäre ebenfalls im Hinblick auf die Produktionsreife empfehlenswert.

Insgesamt zeigt die vorliegende Arbeit, dass die \gls{btp} ein leistungsfähiges Fundament für die Entwicklung integrierter Geschäftsprozessanwendungen bietet \citep{newman2021microservices}. Die im Ausblick skizzierten Erweiterungen verdeutlichen zugleich das Potenzial, den vorgestellten Prototyp schrittweise zu einer produktionsreifen Lösung auszubauen, die den Anforderungen realer Unternehmensumgebungen gerecht wird.


\pagenumbering{Roman}
\setcounter{page}{\value{savepage}}

% Literaturverzeichnis
\cleardoublepage
\phantomsection
\addcontentsline{toc}{section}{Literatur}
\bibliographystyle{dinat}
\bibliography{Bibliothek}

% Übersicht verwendeter Hilfsmittel
\cleardoublepage
\phantomsection
\addcontentsline{toc}{section}{Übersicht verwendeter Hilfsmittel}
\section*{Übersicht verwendeter Hilfsmittel}
\markboth{Übersicht verwendeter Hilfsmittel}{Übersicht verwendeter Hilfsmittel}

Im Rahmen der Erstellung dieser Arbeit wurden neben klassischen Hilfsmitteln (Literaturquellen, Schreib- und Satzprogrammen) auch generative IT-/KI-Werkzeuge unterstützend eingesetzt.

\medskip
Die verwendeten Werkzeuge im Einzelnen:

\begin{itemize}
  \item \textbf{Claude Code (\url{https://claude.ai}):}
  Unterstützung bei der Implementierung des CAP-Services, der CDS-Modelle und der Workflow-Integration. Das Werkzeug wurde zur Codegenerierung, Fehleranalyse und iterativen Verbesserung der Anwendung eingesetzt. Alle generierten Codeartefakte wurden geprüft, angepasst und im Kontext der Anwendungsarchitektur validiert. Darüber hinaus wurde Claude Code zur sprachlichen Unterstützung bei der Erstellung dieser Hausarbeit verwendet.

  \item \textbf{SAP Business Application Studio:}
  Cloudbasierte Entwicklungsumgebung für die Erstellung und das Testen der CAP-Anwendung sowie der Fiori-Elements-Oberfläche.

  \item \textbf{SAP Build Process Automation:}
  Low-Code-Werkzeug zur Modellierung des Genehmigungsworkflows mit Formularen, Bedingungen und API-Aktionen.

  \item \textbf{Overleaf / \LaTeX{}:}
  Satz und Formatierung der vorliegenden Hausarbeit.
\end{itemize}


% Anhang
\cleardoublepage
\phantomsection
\addcontentsline{toc}{section}{Anhang}
\section*{Anhang}
\markboth{Anhang}{Anhang}

\subsection*{A\quad Domänenmodell (schema.cds)}
\begin{lstlisting}[caption={Vollständiges CDS-Schema},label={lst:anhang-schema}]
namespace purchase.approval;
using { cuid, managed, Currency } from '@sap/cds/common';

type RequisitionStatus : String(20) enum {
  Draft; Submitted; Approved; Rejected;
}

entity PurchaseRequisitions : cuid, managed {
  description       : String(256)  @mandatory;
  requester         : String(128)  @mandatory;
  costCenter        : String(10)   @mandatory;
  totalAmount       : Decimal(15,2);
  currency          : Currency;
  status            : RequisitionStatus default 'Draft';
  workflowId        : String(256);
  statusCriticality : Integer default 0 @UI.Hidden;
  items             : Composition of many Items
                        on items.parent = $self;
}

entity Items : cuid {
  parent        : Association to PurchaseRequisitions;
  materialNo    : String(40)   @mandatory;
  description   : String(256);
  quantity      : Decimal(13,3) @mandatory;
  unitPrice     : Decimal(15,2) @mandatory;
  currency      : Currency;
  netAmount     : Decimal(15,2);
}
\end{lstlisting}

\subsection*{B\quad Servicedefinition (purchase-service.cds)}
\begin{lstlisting}[caption={CDS-Servicedefinition mit Draft-Support und Aktionen},label={lst:anhang-service}]
using { purchase.approval as db } from '../db/schema';

service PurchaseService @(path: '/purchase') {
  @odata.draft.enabled
  entity PurchaseRequisitions
    as projection on db.PurchaseRequisitions actions {
      action submit() returns PurchaseRequisitions;
  };
  entity Items as projection on db.Items;

  action workflowCallback (
    workflowId : String,
    status     : String,
    comment    : String
  ) returns PurchaseRequisitions;

  function getSpaInboxUrl() returns String;
}
\end{lstlisting}

\subsection*{C\quad OpenAPI-Spezifikation (Callback-Action)}
\begin{lstlisting}[caption={OpenAPI-Spezifikation für den Workflow-Callback},label={lst:anhang-openapi}]
{
  "openapi": "3.0.0",
  "info": {
    "title": "Purchase Approval Callback",
    "version": "1.0.0"
  },
  "paths": {
    "/purchase/workflowCallback": {
      "post": {
        "operationId": "workflowCallback",
        "requestBody": {
          "content": {
            "application/json": {
              "schema": {
                "properties": {
                  "workflowId": { "type": "string" },
                  "status": {
                    "type": "string",
                    "enum": ["APPROVED", "REJECTED"]
                  },
                  "comment": { "type": "string" }
                }
              }
            }
          }
        }
      }
    }
  }
}
\end{lstlisting}

\subsection*{E\quad Submit-Aktion mit Workflow-Start (purchase-service.js -- Auszug)}
\begin{lstlisting}[caption={Submit-Aktion: Validierung, Statuswechsel und Workflow-Start},label={lst:anhang-submit}]
this.on('submit', PurchaseRequisitions, async req => {
  const id = req.params[0]?.ID || req.params[0];

  const pr = await SELECT.one.from(PurchaseRequisitions)
    .where({ ID: id })
    .columns(pr => { pr('*'); pr.items(i => i('*')); pr.attachments(a => a('*')); });

  if (!pr) return req.error(404, `Bestellanforderung ${id} nicht gefunden`);
  if (pr.status === 'Submitted')
    return req.error(409, 'Bereits eingereicht');
  if (!pr.items || pr.items.length === 0)
    return req.error(422, 'Mindestens eine Position erforderlich');
  if (pr.totalAmount <= 0)
    return req.error(422, 'Gesamtbetrag muss groesser als 0 sein');

  await UPDATE(PurchaseRequisitions).set({ status: 'Submitted' }).where({ ID: id });

  try {
    const workflowInstanceId = await workflowClient.startApprovalWorkflow(pr);
    await UPDATE(PurchaseRequisitions)
      .set({ workflowId: workflowInstanceId }).where({ ID: id });
  } catch (err) {
    await UPDATE(PurchaseRequisitions).set({ status: 'Draft' }).where({ ID: id });
    return req.error(502, `Workflow konnte nicht gestartet werden: ${err.message}`);
  }
});

this.on('workflowCallback', async req => {
  const { workflowId, status, comment, approverName, decidedAt } = req.data;

  if (!workflowId) return req.error(400, 'workflowId ist erforderlich');
  if (!['APPROVED', 'REJECTED'].includes(status))
    return req.error(400, `Ungueltiger Status "${status}"`);

  const pr = await SELECT.one.from(PurchaseRequisitions).where({ workflowId });
  if (!pr) return req.error(404, `Kein Datensatz fuer Workflow ${workflowId}`);
  if (pr.status !== 'Submitted') return req.error(409, 'Datensatz nicht im Status Submitted');

  const callbackStatus = status === 'APPROVED' ? 'Approved' : 'Rejected';
  await UPDATE(PurchaseRequisitions).set({
    status: callbackStatus,
    approverComment: comment || null,
    approvedBy: approverName || null,
    decidedAt: new Date(decidedAt || Date.now()).toISOString(),
    ...(callbackStatus === 'Rejected' ? { workflowId: null } : {})
  }).where({ ID: pr.ID });
});
\end{lstlisting}

\subsection*{F\quad WorkflowClient (workflow-client.js -- Auszug)}
\begin{lstlisting}[caption={WorkflowClient: OAuth2-Authentifizierung und Workflow-Start via SPA-API},label={lst:anhang-workflowclient}]
class WorkflowClient {

  async _getToken() {
    const response = await fetch(process.env.SPA_TOKEN_URL, {
      method: 'POST',
      headers: {
        'Content-Type': 'application/x-www-form-urlencoded',
        'Authorization': 'Basic ' + Buffer.from(
          `${process.env.SPA_CLIENT_ID}:${process.env.SPA_CLIENT_SECRET}`
        ).toString('base64')
      },
      body: 'grant_type=client_credentials'
    });
    const data = await response.json();
    return data.access_token;
  }

  async startApprovalWorkflow(requisition) {
    const token = await this._getToken();
    const workflowContext = {
      requisitionId: requisition.ID,
      description:   requisition.description,
      requester:     requisition.requester,
      costCenter:    requisition.costCenter,
      totalAmount:   parseFloat(requisition.totalAmount),
      currency:      requisition.currency_code || 'EUR',
      items: requisition.items.map(item => ({
        materialNo:  item.materialNo,
        description: item.description,
        quantity:    String(item.quantity),
        unitPrice:   String(item.unitPrice),
        netAmount:   String(item.netAmount)
      })),
      detailLink: `${process.env.APP_URL}#/PurchaseRequisitions(ID=${requisition.ID},IsActiveEntity=true)`,
      callbackUrl: process.env.CAP_CALLBACK_URL
    };

    const result = await fetch(
      `${process.env.SPA_API_URL}/workflow/rest/v1/workflow-instances`,
      {
        method: 'POST',
        headers: {
          'Authorization': `Bearer ${token}`,
          'Content-Type': 'application/json'
        },
        body: JSON.stringify({
          definitionId: process.env.SPA_WORKFLOW_DEFINITION_ID,
          context: workflowContext
        })
      }
    );
    const data = await result.json();
    return data.id;
  }
}
\end{lstlisting}

\subsection*{D\quad MTA-Deskriptor (mta.yaml -- Auszug)}
\begin{lstlisting}[caption={MTA-Deskriptor mit Modulen und Ressourcen},label={lst:anhang-mta}]
ID: purchase-approval
version: 1.0.0
modules:
  - name: purchase-approval-srv
    type: nodejs
    path: gen/srv
    requires:
      - name: purchase-approval-db
      - name: purchase-approval-auth
      - name: purchase-approval-destination
    provides:
      - name: srv-api
        properties:
          srv-url: ${default-url}

  - name: purchase-approval-db-deployer
    type: hdb
    path: gen/db
    requires:
      - name: purchase-approval-db

  - name: purchase-approval-app
    type: approuter.nodejs
    path: app
    requires:
      - name: srv-api
      - name: purchase-approval-auth

resources:
  - name: purchase-approval-db
    type: com.sap.xs.hdi-container
  - name: purchase-approval-auth
    type: org.cloudfoundry.managed-service
    parameters:
      service: xsuaa
      service-plan: application
      path: ./xs-security.json
  - name: purchase-approval-destination
    type: org.cloudfoundry.managed-service
    parameters:
      service: destination
      service-plan: lite
\end{lstlisting}


\end{document}
